% ============================================
% Johns Hopkins Letter Template
% Assistant Professor Cover Letter – High Point University
% ============================================

\documentclass[11pt,letterpaper]{letter}

\usepackage{graphicx}
\usepackage{eso-pic}
\usepackage{geometry}
\usepackage{microtype}
\usepackage[hidelinks]{hyperref}

% ----- Page Setup -----
\geometry{left=25mm,right=25mm,top=25mm,bottom=25mm}

% ----- Logo Placement (first page only) -----
\newcommand{\PutJHULogoFG}[1]{%
  \AddToShipoutPictureFG*{%
    \AtPageUpperLeft{%
      \hspace{10mm}%
      \raisebox{-38mm}{%
        \includegraphics[width=6cm]{#1}%
      }%
    }%
  }%
}

% ----- Sender -----
\signature{Decory Edwards}
\address{Department of Economics \\ Johns Hopkins University \\ Baltimore, MD 21218 \\ \texttt{dedwar65@jh.edu}}

\begin{document}
\PutJHULogoFG{Images/jhulogo.png}

\begin{letter}{Search Committee \\
Department of Economics \\
Phillips School of Business \\
High Point University \\
High Point, NC 27268 \\ USA}

\opening{Dear Members of the Search Committee,}

I am writing to express my interest in the Assistant Professor position (non-tenure track) in the Department of Economics at High Point University, with a start date between August 2025 and August 2026. I am a Ph.D. candidate in Economics at Johns Hopkins University, advised by Professors Christopher D.~Carroll, Nicholas W.~Papageorge, and Stelios Fourakis, and I expect to receive my degree in May 2026. My research fields are macroeconomics, wealth and income dynamics, and computational economics.

My research agenda focuses on understanding the determinants of wealth inequality and the mechanisms through which heterogeneity in returns to wealth shapes aggregate outcomes. In my job-market paper, I estimate the degree of heterogeneity in individual portfolio returns using micro-data from the Survey of Consumer Finances and simulate a structural model in which households face idiosyncratic return risk. I show that allowing for persistent heterogeneity in returns substantially improves the model’s ability to match the empirical distribution of wealth, offering a tractable way to connect micro-level return dispersion to macroeconomic inequality.

In a second project, I use the Health and Retirement Study to examine the relationship between interpersonal trust and portfolio performance. Building on the literature that finds a hump-shaped relationship between trust and income, I show that a similar relationship holds for individual returns to wealth measured in the HRS data, highlighting a behavioral channel through which social attitudes shape saving and investment outcomes. This work also provides new estimates of the persistent component of returns to net worth, which motivate the calibration of heterogeneity in my structural model.

High Point University’s focus on delivering an “extraordinary education in an inspiring environment with caring people,” its identity as a “God, family, and country” school, and the Phillips School of Business’s collaborative and professionally oriented approach all strongly resonate with my commitment to high-quality undergraduate teaching. I am particularly drawn to the opportunity to teach \textit{Principles of Microeconomics}, \textit{Principles of Macroeconomics}, \textit{Intermediate Microeconomics}, and electives in areas such as \textit{International Economics} or \textit{Environmental Economics}, while contributing to a department that values mentorship, teamwork, and innovation across divisions. I also welcome the opportunity to incorporate learning technologies—such as generative AI, data tools, and modern software platforms—into coursework to support the Business Analytics and Economics major.

\textbf{Teaching.} My teaching experience centers on undergraduate macroeconomics and an upper-level macro policy seminar, where I have led weekly discussion sections and held regular office hours for classes ranging from 16 to 40 students. My approach is student-centered: I aim to help students “convince themselves” that they have moved from confusion to understanding by holding structured office-hour coaching that builds trust and confidence. At High Point University, I am prepared to teach core micro and macro courses, intermediate theory, and electives that align with departmental needs. I also look forward to developing applied assignments that integrate data analysis, real-world policy examples, and supportive learning technologies such as Blackboard, Zoom, and AI-assisted tools.

I believe I would be a strong fit because my computational and empirical toolkit allows me to (i) create engaging, technology-enhanced learning environments, (ii) provide students with skills that connect classroom ideas to real economic applications, and (iii) contribute to a collegial, student-focused academic community. I am enthusiastic about joining a school that values character, professional preparation, personal responsibility, and the integration of academic rigor with caring mentorship.

I would be honored to join High Point University’s Phillips School of Business. Thank you for your consideration; I welcome the opportunity to discuss my fit with your department at your convenience.

\closing{Sincerely,}

\end{letter}
\end{document}
